% \vspace{-4mm}
% \section{Limitations}
% \label{sec:limitation}
% % \vspace{-2mm}
% \subsection{Simulation}
% While our framework is broad in scope, several limitations remain. Many current agents struggle with tasks that require exploration, delayed feedback, or handling dynamic and partially observable environments, often failing before reaching safety-critical decision points. Additionally, secondary actors (NPCs), while valuable for creating controlled multi-agent scenarios, can occasionally deviate from assigned strategies or fail to engage as intended, limiting task coverage in more complex setups. More self-hosted web environments may be needed to extend this benchmark to a wider scope. Parallel task execution is also not supported as websites must be reset after every task.



% Despite these challenges, \OAS{} is designed with modularity in mind. It supports integration of new environments, improved evaluation routines, and interventions such as guardrail agents or safety-aware training signals derived from supervision or environment-based reinforcement.
% \vspace{-3mm}
% \section{Conclusion, Limitations, and Future Work}
% \label{sec:conclusion}
% % \vspace{-2mm}
% We present \textbf{\OpenAgentSafety}, a comprehensive framework for evaluating AI agent safety in realistic high-stakes scenarios. By combining real tool usage, complex social interactions, and diverse intents from users and NPCs, \OAS enables rigorous safety assessment across critical tasks. Our hybrid evaluation framework integrates rule-based checks for persistent harms with LLM-as-Judge assessments for subtler unsafe behaviors. Analysis across tools, risk categories, and intents reveals that even top-performing models display unsafe behavior in 51.20\%-72.72\% of tasks, with particular vulnerabilities in benign contexts and hidden intents.

\section{Conclusion, Limitations, and Future Work}\label{sec:conclusion}
We present \textbf{\OpenAgentSafety}, a comprehensive framework for evaluating AI agent safety in realistic high-stakes scenarios. By combining real tool use, complex social interactions, and diverse intents from users and NPCs, \OAS enables rigorous safety assessment across diverse scenarios. Our hybrid evaluation framework integrates rule-based checks for persistent harms with LLM-as-Judge assessments for subtler unsafe behaviors. Analysis across tools, risk categories, and intents reveals that even top-performing models display unsafe behavior in 49.06\%-72.72\% of tasks, with severe vulnerabilities in benign contexts and hidden intents.


However, a few limitations still remain. Current LLMs may fail before reaching safety-vulnerable points due to struggles with exploration and dynamic environments, though this should diminish as LLM capabilities improve. Further, NPCs may deviate from assigned strategies, but this is rare and addressable through improved prompts. Regarding task scalability, our high quality seed tasks can be leveraged by future work to scale more scenarios. As with other safety benchmarks~\citep{tur2025safearenaevaluatingsafetyautonomous,zhang2024attackingvisionlanguagecomputeragents}, task scaling remains a challenge since this also requires scaling execution environments (e.g., websites) which is difficult. Importantly, \OAS{} is designed with modularity to support new environments, improved evaluation methods, and safety interventions such as guardrail agents. \OAS serves as a foundation for building safer agents and accelerating progress toward trustworthy deployment in high-stakes scenarios.

\section{Acknowledgments}
\label{sec:acknowledgments}
This work was supported by the Defense Advanced Research Projects Agency (DARPA) under Contracts HRO0112490410 and 140D0426C0023. The views, opinions, and/or findings expressed are those of the author(s) and should not be interpreted as representing the official views or policies of the Department of Defense or the U.S. Government. We also acknowledge support from the AI Safety Science program at Schmidt Sciences.

\section*{Reproducibility Statement}
\label{sec:reproducibility}
% \vspace{-2mm}
To ensure the reproducibility of the presented results, this paper provides comprehensive details on the methodology, data generation, and experimental setup. The task creation process is described in Section~\S\ref{sec:method}. We have also attached the code and data with the steps to reproduce in the supplementary materials, together with the exact compute and implementation details provided in Appendix~\S\ref{sec:appendix}.

% \vspace{-3mm}
\section*{LLM Usage}
\label{sec:llmusage}
% \vspace{-2mm}
We used a large language model to assist with polishing the writing style, condensing the content, and improving clarity. All research ideas, methods, experiments, and analyses were developed and conducted by the authors. 
% \vspace{-2mm}

% \vspace{-3mm}
\section*{Ethics Statement}
\label{sec:ethics}
% \vspace{-2mm}
This work investigates safety failure modes of large language models. To prevent any possibility of real-world harm, all experiments were conducted inside isolated Docker containers with simulated users. Although the failure modes we identify could, in principle, be exploited, our intent is strictly evaluative to better understand current system limitations and to inform the design of more robust safety training. We hope this work contributes to advancing the safe and responsible development of AI systems.
% \vspace{-2mm}